% Generated draft: requires mathematical and source review.
\documentclass[11pt]{article}
\usepackage[margin=25mm]{geometry}
\usepackage{fontspec}
\setmainfont{Latin Modern Roman}
\usepackage{amsmath,amssymb,mathtools,unicode-math}
\setmathfont{Latin Modern Math}
\usepackage{enumitem,needspace,xurl,hyperref}
\setlength{\parindent}{0pt}
\setlength{\parskip}{5pt}
\setlength{\emergencystretch}{3em}
\newcommand{\N}{\mathbb{N}}
\newcommand{\Z}{\mathbb{Z}}
\newcommand{\Q}{\mathbb{Q}}
\newcommand{\R}{\mathbb{R}}
\newcommand{\C}{\mathbb{C}}
\newcommand{\F}{\mathbb{F}}
\newcommand{\A}{\mathbb{A}}
\newcommand{\PP}{\mathbb{P}}
\newcommand{\E}{\mathbb{E}}
\newcommand{\eps}{\varepsilon}
\newcommand{\dd}{\,\mathrm{d}}
\begin{document}
\section*{UnsolvedMath ID: 30001684}
\textbf{Source problem code:} catalogue entry
\subsection{Definitions and mathematical statement}
Work in a ground model \(V\) of \(\mathrm{ZFC}\). The generalized continuum hypothesis (GCH) is the assertion that \(2^\mu=\mu^+\) for every infinite cardinal \(\mu\). A cardinal \(\kappa\) is strongly inaccessible if it is uncountable, regular, and satisfies \(2^\mu<\kappa\) for every cardinal \(\mu<\kappa\).

Write
\[
[\kappa^+]^{<\kappa}=\{x\subseteq\kappa^+:|x|<\kappa\}.
\]
A subset of this set is club if it is unbounded under inclusion and closed under unions of increasing chains of nonzero length less than \(\kappa\). A subset is stationary if it meets every such club. For a set of ordinals \(x\), let \(\operatorname{otp}(x)\) denote its order type in the inherited ordering. Define
\[
\mathrm{SC}(\kappa)\quad\Longleftrightarrow\quad
\kappa\text{ is strongly inaccessible and }
S_\kappa=\{x\in[\kappa^+]^{<\kappa}:\operatorname{otp}(x)\text{ is a cardinal}\}
\text{ is stationary}.
\]
Here a cardinal is identified with its initial ordinal, and these definitions are evaluated in \(V\).

For an infinite cardinal \(\lambda\), use the following convention for \(\square_{\lambda,<\omega}\): there is a sequence
\[
\langle\mathcal C_\alpha:\alpha<\lambda^+,\ \alpha\text{ a nonzero limit ordinal}\rangle
\]
such that each \(\mathcal C_\alpha\) is a nonempty finite collection of closed unbounded subsets of \(\alpha\), each member has order type at most \(\lambda\), and whenever \(C\in\mathcal C_\alpha\) and \(\beta<\alpha\) is a limit point of \(C\), one has \(C\cap\beta\in\mathcal C_\beta\). There is no club \(D\subseteq\lambda^+\) such that \(D\cap\alpha\in\mathcal C_\alpha\) at every limit point \(\alpha<\lambda^+\) of \(D\). A limit point \(\beta\) of a set \(C\) satisfies \(\beta>0\) and \(\sup(C\cap\beta)=\beta\). Closedness means containing all limit points below the ambient ordinal, and unboundedness means being cofinal in that ordinal.

\textbf{Extracted mathematical question.} Given \(\kappa\) with \(V\models\mathrm{GCH}+\mathrm{SC}(\kappa)\), can one obtain a forcing extension \(V[G]\) of the intended Prikry kind such that
\[
V[G]\models\neg\square_{\aleph_\omega,<\omega},
\]
without adding the hypothesis that \(\kappa\) carries a normal measure? The cardinals in this displayed conclusion are evaluated in the extension. A normal measure on \(\kappa\) means a nonprincipal \(\kappa\)-complete ultrafilter on \(\kappa\) such that every regressive function on a measure-one set is constant on a measure-one subset. Omitting this hypothesis does not require assuming that no normal measure exists.

\textbf{Unresolved scope.} The original question is whether the assumption that \(\kappa\) is measurable can be eliminated from Theorem 1. The supplied record does not reproduce that theorem or specify the intended forcing. Ordinary Prikry forcing uses a normal measure, so the phrase ``Prikry extension'' requires clarification when that hypothesis is removed. The record also does not specify the relationship between the ground-model \(\kappa\) and the extension's \(\aleph_\omega\), or any additional preservation requirements. Thus the displayed question records the supplied extracted target; a fully specified replacement for Theorem 1 cannot be recovered from this record alone.
\subsection{Short English statement}
Assuming GCH and the stated stationary-set principle \(\mathrm{SC}(\kappa)\), can the intended Prikry extension make finite-width square fail at \(\aleph_\omega\) without an additional measurability assumption on \(\kappa\)? The referenced theorem and the precise meaning of the extension are missing, so its full hypotheses and conclusion require review.
\subsection{Sources}
\begin{itemize}
\item Supplied local UnsolvedMath registry record, permanent ID \(30001684\), source problem code \texttt{catalogue entry}, title: Measurability Assumption in the Prikry Extension Theorem. The record assigns the problem to 2011 and supplies no individual proposer.
\item \emph{Set Theory} (2011), Oberwolfach Reports, pp.\ 85--140, DOI \texttt{10.4171/owr/2011/02}: \url{https://doi.org/10.4171/owr/2011/02}. This is the supplied report citation and the literature-source link labelled ``Oberwolfach problem statement.'' The original question refers to Theorem 1, whose text is absent from the supplied record.
\item Supplied literature source labelled ``Schimmerling and Zeman, square in core models'': \url{https://doi.org/10.2307/2687750}.
\item Supplied literature source labelled ``consistency strength of failure of square'': \url{https://www.math.uci.edu/~mzeman/RESEARCH/ffs-0-120719.pdf}.
\item Literature assessment dated exactly \texttt{2026-08-21}: the supplied curation labels the problem open and reports that the cited Prikry-extension theorem uses measurability of \(\kappa\), with no elimination of that hypothesis found in the literature checked. This is supplied dated curation, not an independently verified literature assessment.
\item The registry labels the statement \texttt{corrected\_verified} and reports that the official report was checked to restore the large-cardinal principle, forcing conclusion, and measurability question. These are supplied verification claims; the linked sources have not been independently consulted for this draft.
\end{itemize}
\end{document}
